\documentclass{article}

\title{ReSTIRSSS Ergebnisse}
\author{Felix Viering}
\begin{document}
\maketitle
\section{Versuch}
Die Szene "Ajax" wurde auf einem Linux-System mit einer AMD Radeon RX 9070 XT
Grafikkarte, welche hardwarebeschleunigtes Raytracing unterstützt, getestet.
Die Einstellung \emph{Max Frames} wurde auf den Wert 2048 eingestellt und es
wurden die Presets \emph{Path Tracing}, \emph{ReSTIRSSS (Hybrid Shift)},
\emph{ReSTIRSSS (Sequential Shift)} verwendet. Es wurden die Zeiten aufgenommen,
die das Rendern eines Frames bei den unterschiedlichen Einstellungen dauert und
die resultierenden Bilder nach dem Rendern der maximalen Anzahl an Frames verglichen.

\section{Ergebnisse}
Die Renderzeiten lassen sich aus Tabelle 1 ableiten. 
\begin{table}
    \caption{Renderzeiten}\label{tab:}
    \begin{center}
        \begin{tabular}[c]{|l|l|l|}
            \hline
            \multicolumn{1}{|c|}{\textbf{Path Tracing}} & 
            \multicolumn{1}{c|}{\textbf{ReSTIRSSS (Hybrid Shift)}} &
            \multicolumn{1}{c|}{\textbf{ReSTIRSSS (Sequential Shift)}}\\
            \hline
            0,839ms & 4,319ms & 4,784ms \\


            
            \hline
        \end{tabular}
    \end{center}
\end{table}

\begin{figure}
    \begin{center}
        \includegraphics[width=0.95\textwidth]{./PathTracing.png.hdr}
    \end{center}
    \caption{}\label{fig:}
\end{figure}



\end{document}
